% AulaTeX - maqueta inicial de construcción descendente
% Contrato futuro: el Agente puede investigar, redactar, evaluar y compilar a partir de este archivo.
\documentclass[12pt]{article}
\usepackage[utf8]{inputenc}
\usepackage[T1]{fontenc}
\usepackage[spanish]{babel}
\usepackage{enumitem}
\usepackage{hyperref}
\usepackage{longtable}

\title{UCNL}
\author{AulaTeX}
\date{\today}

\begin{document}
\maketitle

\noindent\textbf{Nodo editorial:} UCNL\\
\textbf{Padre editorial:} interinstitucional\\
\textbf{Modo de generación:} reforzar\\
\textbf{Destino:} UCNL\\
\textbf{Contrato futuro:} ready-for-agent

\section*{Ingesta base}
\begin{itemize}[leftmargin=*]
  \item Texto libre proporcionado: no
  \item Documento de apoyo proporcionado: D:\\textbackslash\\{\\}Descargas\\textbackslash\\{\\}Modelo Educativo UC.pdf
\end{itemize}

\section*{Objetivo}
\begin{itemize}[leftmargin=*]
  \item Proveer un esqueleto reusable para construir entregables UCNL sin romper reglas editoriales ni tecnicas.
  \item Servir como esqueleto reusable para entregables UCNL.
  \item Completar la maqueta inicial de UCNL.
\end{itemize}

\section*{Competencias}
\begin{itemize}[leftmargin=*]
  \item Redaccion academica argumentada.
  \item Sintesis y analisis de fuentes.
  \item Gestion de referencias con BibTeX/natbib.
  \item Produccion documental en LaTeX.
  \item Redaccion academica.
  \item Analisis de fuentes.
  \item Gestion bibliografica en LaTeX.
  \item Presentacion formal de resultados.
\end{itemize}

\section*{Resultados esperados}
\begin{itemize}[leftmargin=*]
  \item Documento compilable en PDF con estructura academica completa.
  \item Citas y referencias coherentes con fuentes usadas.
  \item Cumplimiento de producto solicitado por consigna.
  \item PDF compilable con estructura academica completa.
  \item Uso correcto de citas y referencias.
  \item Cumplimiento del producto solicitado.
\end{itemize}

\section*{Estructura sugerida}
\begin{itemize}[leftmargin=*]
  \item Portada institucional.
  \item Resumen breve.
  \item Introduccion y contexto.
  \item Desarrollo por criterios/rubrica.
  \item Producto solicitado (tabla, esquema, analisis, etc.).
  \item Conclusion.
  \item Referencias bibliograficas.
  \item Resumen.
  \item Introduccion.
  \item Desarrollo alineado a criterios.
  \item Producto solicitado.
  \item Referencias.
\end{itemize}

\section*{Criterios de evaluación}
\begin{itemize}[leftmargin=*]
  \item Pertinencia y profundidad del contenido.
  \item Organizacion y coherencia argumentativa.
  \item Correccion formal de citacion y referencias.
  \item Calidad de formato y compilacion final.
  \item Pertinencia del contenido.
  \item Coherencia argumentativa.
  \item Correccion formal.
  \item Calidad de compilacion.
\end{itemize}

\section*{Bibliografía requerida}
\begin{itemize}[leftmargin=*]
  \item bibliografia-ucnl.bib como base.
  \item Fuentes de planeacion y documentos institucionales aplicables.
  \item bibliografia-ucnl.bib.
  \item Fuentes institucionales aplicables.
\end{itemize}

\section*{Marcadores de investigación}
\begin{itemize}[leftmargin=*]
  \item Pendiente: extraer lineamientos verificables del Modelo Educativo UC.
  \item Pendiente: identificar rubrica especifica de materia/actividad.
  \item Pendiente: definir corpus minimo de fuentes por entrega.
  \item Extraer lineamientos verificables del Modelo Educativo UC.
  \item Definir rubrica especifica.
  \item Delimitar corpus minimo de fuentes.
  \item Definir fuentes, citas y vacíos de contenido antes de redactar.
\end{itemize}

\section*{Indicaciones editoriales por entregable}

\subsection*{Plantilla}
\begin{itemize}[leftmargin=*]
  \item Archivo base institucional derivable desde UCNL/reporte-ucnl.tex.
  \item Bloques obligatorios: metadatos, resumen, desarrollo, conclusion, bibliografia.
  \item Variables de portada y marca de agua configurables sin tocar estructura troncal.
  \item Plantilla institucional UCNL como base.
  \item Bloques obligatorios predefinidos.
  \item Variables configurables sin alterar estructura.
  \item Definir plantilla base con reglas editoriales, assets y referencias de estilo antes de redactar.
  \item Prepara una raíz institucional reutilizable con carpetas por carrera, carpeta assets y referencias claras a plantillas LaTeX compartidas.
  \item Compilar mediante scripts institucionales pasando solo el .tex de entrada.
  \item Resolver plantillas por TEXINPUTS y bibliografia por BIBINPUTS definidos en .latexmkrc.
  \item Usar natbib para citacion autor-anio (\\textbackslash\\{\\}citet, \\textbackslash\\{\\}citep) y mantener compatibilidad APA 7 en salida.
  \item Conservar raiz UCNL con: bibliografia institucional, referencias institucionales, assets y subcarpetas academicas.
  \item Cada materia/carrera debe incluir COMPILACION.md con comando exacto, contrato de compilacion y bib esperado.
  \item Separar plantilla maestra (institucional) de entregables por actividad (nivel materia/semana).
  \item Usa la ingesta proporcionada como restricción editorial inicial.
\end{itemize}

\subsection*{Actividad}
\begin{itemize}[leftmargin=*]
  \item Plantilla hija por materia/semana con secciones guiadas y placeholders de evidencia.
  \item Checklist embebido en comentarios para cumplimiento de consigna y rubrica.
  \item Sin texto final redactado; solo andamiaje editorial.
  \item Plantilla hija con secciones guiadas.
  \item Placeholders y checklist en comentarios.
  \item Sin texto final redactado.
  \item Desarrollar la actividad respetando objetivo, criterios de evaluación y marcadores de investigación.
  \item Proveer un esqueleto reusable para construir entregables UCNL sin romper reglas editoriales ni tecnicas.
  \item Servir como esqueleto reusable para entregables UCNL.
  \item Pertinencia y profundidad del contenido.
  \item Organizacion y coherencia argumentativa.
  \item Correccion formal de citacion y referencias.
  \item Validar lineamientos formales del Modelo Educativo UC antes de redactar actividades.
  \item Identificar criterios de evaluacion explicitos por materia y mapearlos a secciones del reporte.
  \item Confirmar formato de entrega solicitado en aula (nomenclatura, extension, producto visual).
  \item Usa la ingesta proporcionada como restricción editorial inicial.
\end{itemize}

\subsection*{Reporte}
\begin{itemize}[leftmargin=*]
  \item Entrada canonica por producto academico con natbib activo y bibliografia enlazada.
  \item Secciones orientadas a analisis y producto solicitado.
  \item Compatible con compilacion por scripts institucionales.
  \item Entrada canonica con natbib activo.
  \item Orientada a analisis y evidencia.
  \item Compatible con scripts institucionales.
  \item Estructurar el reporte con bibliografía requerida y cierre verificable.
  \item UCNL/reporte-ucnl.tex (plantilla maestra institucional).
  \item UCNL/presentacion-ucnl.tex (entrada canonica para diapositivas).
  \item UCNL/bibliografia-ucnl.bib (fuente bibliografica institucional).
  \item Redaccion academica formal, precisa y sin copiar instrucciones de la actividad.
  \item Priorizar sintesis, analisis y postura fundamentada sobre transcripcion.
  \item Mantener consistencia terminologica entre portada, cuerpo, figuras y referencias.
  \item Centralizar fuentes institucionales en bibliografia-ucnl.bib y complementar en bib de materia cuando aplique.
  \item No inventar referencias; registrar solo fuentes verificables usadas en el texto.
  \item Toda cita en cuerpo debe corresponder a una entrada valida en .bib.
  \item Cumple exactamente el producto solicitado (ensayo, cuadro, mapa, reporte, presentacion, etc.).
  \item Incluye citas en texto y referencias bibliograficas trazables en .bib institucional o de materia.
  \item Compila sin errores bloqueantes y genera PDF en carpeta objetivo.
\end{itemize}

\subsection*{Presentación}
\begin{itemize}[leftmargin=*]
  \item Entrada canonica UCNL/presentacion-ucnl.tex para sintesis visual del reporte.
  \item Estructura minima: objetivo, hallazgos, evidencia, cierre, referencias.
  \item Consistencia terminologica con el reporte escrito.
  \item Entrada canonica para sintesis visual.
  \item Estructura minima estandar.
  \item Consistencia con el reporte escrito.
  \item Preparar la presentación con síntesis visual, continuidad editorial y evidencias clave.
  \item Documento compilable en PDF con estructura academica completa.
  \item Citas y referencias coherentes con fuentes usadas.
  \item Cumplimiento de producto solicitado por consigna.
  \item Redaccion academica formal, precisa y sin copiar instrucciones de la actividad.
  \item Priorizar sintesis, analisis y postura fundamentada sobre transcripcion.
  \item Nivel institucion: normas de identidad, estructura, estilo, calidad y compilacion.
  \item No incluye redaccion completa de actividades ni investigacion de fondo en este ciclo.
  \item Sintetiza visualmente la memoria editorial sin perder continuidad con plantilla, actividad y reporte.
\end{itemize}

\end{document}
